%%
% Copyright (c) 2017 - 2024, Pascal Wagler;
% Copyright (c) 2014 - 2024, John MacFarlane
%
% All rights reserved.
%
% Redistribution and use in source and binary forms, with or without
% modification, are permitted provided that the following conditions
% are met:
%
% - Redistributions of source code must retain the above copyright
% notice, this list of conditions and the following disclaimer.
%
% - Redistributions in binary form must reproduce the above copyright
% notice, this list of conditions and the following disclaimer in the
% documentation and/or other materials provided with the distribution.
%
% - Neither the name of John MacFarlane nor the names of other
% contributors may be used to endorse or promote products derived
% from this software without specific prior written permission.
%
% THIS SOFTWARE IS PROVIDED BY THE COPYRIGHT HOLDERS AND CONTRIBUTORS
% "AS IS" AND ANY EXPRESS OR IMPLIED WARRANTIES, INCLUDING, BUT NOT
% LIMITED TO, THE IMPLIED WARRANTIES OF MERCHANTABILITY AND FITNESS
% FOR A PARTICULAR PURPOSE ARE DISCLAIMED. IN NO EVENT SHALL THE
% COPYRIGHT OWNER OR CONTRIBUTORS BE LIABLE FOR ANY DIRECT, INDIRECT,
% INCIDENTAL, SPECIAL, EXEMPLARY, OR CONSEQUENTIAL DAMAGES (INCLUDING,
% BUT NOT LIMITED TO, PROCUREMENT OF SUBSTITUTE GOODS OR SERVICES;
% LOSS OF USE, DATA, OR PROFITS; OR BUSINESS INTERRUPTION) HOWEVER
% CAUSED AND ON ANY THEORY OF LIABILITY, WHETHER IN CONTRACT, STRICT
% LIABILITY, OR TORT (INCLUDING NEGLIGENCE OR OTHERWISE) ARISING IN
% ANY WAY OUT OF THE USE OF THIS SOFTWARE, EVEN IF ADVISED OF THE
% POSSIBILITY OF SUCH DAMAGE.
%%

%%
% This is the Eisvogel pandoc LaTeX template.
%
% For usage information and examples visit the official GitHub page:
% https://github.com/Wandmalfarbe/pandoc-latex-template
%%

% Options for packages loaded elsewhere
\PassOptionsToPackage{unicode}{hyperref}
\PassOptionsToPackage{hyphens}{url}
\PassOptionsToPackage{dvipsnames,svgnames,x11names,table}{xcolor}
%
\documentclass[
  paper=a4,
  oneside  ,captions=tableheading
]{scrartcl}
\usepackage{amsmath,amssymb}
% Use setspace anyway because we change the default line spacing.
% The spacing is changed early to affect the titlepage and the TOC.
\usepackage{setspace}
\setstretch{1.2}
\usepackage{iftex}
\ifPDFTeX
  \usepackage[T1]{fontenc}
  \usepackage[utf8]{inputenc}
  \usepackage{textcomp} % provide euro and other symbols
\else % if luatex or xetex
  \usepackage{unicode-math} % this also loads fontspec
  \defaultfontfeatures{Scale=MatchLowercase}
  \defaultfontfeatures[\rmfamily]{Ligatures=TeX,Scale=1}
\fi
\usepackage{lmodern}
\ifPDFTeX\else
  % xetex/luatex font selection
\fi
% Use upquote if available, for straight quotes in verbatim environments
\IfFileExists{upquote.sty}{\usepackage{upquote}}{}
\IfFileExists{microtype.sty}{% use microtype if available
  \usepackage[]{microtype}
  \UseMicrotypeSet[protrusion]{basicmath} % disable protrusion for tt fonts
}{}
\makeatletter
\@ifundefined{KOMAClassName}{% if non-KOMA class
  \IfFileExists{parskip.sty}{%
    \usepackage{parskip}
  }{% else
    \setlength{\parindent}{0pt}
    \setlength{\parskip}{6pt plus 2pt minus 1pt}}
}{% if KOMA class
  \KOMAoptions{parskip=half}}
\makeatother
\usepackage{xcolor}
\definecolor{default-linkcolor}{HTML}{A50000}
\definecolor{default-filecolor}{HTML}{A50000}
\definecolor{default-citecolor}{HTML}{4077C0}
\definecolor{default-urlcolor}{HTML}{4077C0}
\usepackage[margin=2.5cm,includehead=true,includefoot=true,centering,]{geometry}
\usepackage[export]{adjustbox}
\usepackage{graphicx}
\usepackage{longtable,booktabs,array}
\usepackage{calc} % for calculating minipage widths
% Correct order of tables after \paragraph or \subparagraph
\usepackage{etoolbox}
\makeatletter
\patchcmd\longtable{\par}{\if@noskipsec\mbox{}\fi\par}{}{}
\makeatother
% Allow footnotes in longtable head/foot
\IfFileExists{footnotehyper.sty}{\usepackage{footnotehyper}}{\usepackage{footnote}}
\makesavenoteenv{longtable}
% add backlinks to footnote references, cf. https://tex.stackexchange.com/questions/302266/make-footnote-clickable-both-ways
\usepackage{footnotebackref}
\usepackage{graphicx}
\makeatletter
\newsavebox\pandoc@box
\newcommand*\pandocbounded[1]{% scales image to fit in text height/width
  \sbox\pandoc@box{#1}%
  \Gscale@div\@tempa{\textheight}{\dimexpr\ht\pandoc@box+\dp\pandoc@box\relax}%
  \Gscale@div\@tempb{\linewidth}{\wd\pandoc@box}%
  \ifdim\@tempb\p@<\@tempa\p@\let\@tempa\@tempb\fi% select the smaller of both
  \ifdim\@tempa\p@<\p@\scalebox{\@tempa}{\usebox\pandoc@box}%
  \else\usebox{\pandoc@box}%
  \fi%
}
% Set default figure placement to htbp
% Make use of float-package and set default placement for figures to H.
% The option H means 'PUT IT HERE' (as  opposed to the standard h option which means 'You may put it here if you like').
\usepackage{float}
\floatplacement{figure}{H}
\makeatother
\setlength{\emergencystretch}{3em} % prevent overfull lines
\providecommand{\tightlist}{%
  \setlength{\itemsep}{0pt}\setlength{\parskip}{0pt}}
\setcounter{secnumdepth}{5}
% Make \paragraph and \subparagraph free-standing
\makeatletter
\ifx\paragraph\undefined\else
  \let\oldparagraph\paragraph
  \renewcommand{\paragraph}{
    \@ifstar
      \xxxParagraphStar
      \xxxParagraphNoStar
  }
  \newcommand{\xxxParagraphStar}[1]{\oldparagraph*{#1}\mbox{}}
  \newcommand{\xxxParagraphNoStar}[1]{\oldparagraph{#1}\mbox{}}
\fi
\ifx\subparagraph\undefined\else
  \let\oldsubparagraph\subparagraph
  \renewcommand{\subparagraph}{
    \@ifstar
      \xxxSubParagraphStar
      \xxxSubParagraphNoStar
  }
  \newcommand{\xxxSubParagraphStar}[1]{\oldsubparagraph*{#1}\mbox{}}
  \newcommand{\xxxSubParagraphNoStar}[1]{\oldsubparagraph{#1}\mbox{}}
\fi
\makeatother
\ifLuaTeX
\usepackage[bidi=basic]{babel}
\else
\usepackage[bidi=default]{babel}
\fi
\babelprovide[main,import]{french}
% get rid of language-specific shorthands (see #6817):
\let\LanguageShortHands\languageshorthands
\def\languageshorthands#1{}
\newcommand{\chapterpage}[1]{
  \newpage
  \vspace*{\fill}
  \begin{center}
  {\Huge\bfseries #1}
  \end{center}
  \vspace*{\fill}
  \newpage
}

\usepackage{colortbl}


% Alternance de couleurs (commence à la ligne 2 pour éviter l'en-tête)
\rowcolors{1}{white}{lightgray!50}

\usepackage[table]{xcolor}
\definecolor{lightgray}{HTML}{F0F0F0}

% Force la coloration à toujours commencer APRES l’en-tête (ligne 2)
\AtBeginEnvironment{tabular}{
  \rowcolors{2}{lightgray}{white}
}

\usepackage{soul}
\sethlcolor{yellow}
\makeatletter
\@ifpackageloaded{caption}{}{\usepackage{caption}}
\AtBeginDocument{%
\ifdefined\contentsname
  \renewcommand*\contentsname{Table des matières}
\else
  \newcommand\contentsname{Table des matières}
\fi
\ifdefined\listfigurename
  \renewcommand*\listfigurename{Liste des Figures}
\else
  \newcommand\listfigurename{Liste des Figures}
\fi
\ifdefined\listtablename
  \renewcommand*\listtablename{Liste des Tables}
\else
  \newcommand\listtablename{Liste des Tables}
\fi
\ifdefined\figurename
  \renewcommand*\figurename{Figure}
\else
  \newcommand\figurename{Figure}
\fi
\ifdefined\tablename
  \renewcommand*\tablename{Table}
\else
  \newcommand\tablename{Table}
\fi
}
\@ifpackageloaded{float}{}{\usepackage{float}}
\floatstyle{ruled}
\@ifundefined{c@chapter}{\newfloat{codelisting}{h}{lop}}{\newfloat{codelisting}{h}{lop}[chapter]}
\floatname{codelisting}{Listing}
\newcommand*\listoflistings{\listof{codelisting}{Liste des Listings}}
\makeatother
\makeatletter
\makeatother
\makeatletter
\@ifpackageloaded{caption}{}{\usepackage{caption}}
\@ifpackageloaded{subcaption}{}{\usepackage{subcaption}}
\makeatother
\usepackage{bookmark}
\IfFileExists{xurl.sty}{\usepackage{xurl}}{} % add URL line breaks if available
\urlstyle{same}
\hypersetup{
  pdfauthor={Hugo Roussaffa},
  pdflang={fr},
  colorlinks=true,
  linkcolor={blue},
  filecolor={default-filecolor},
  citecolor={default-citecolor},
  urlcolor={default-urlcolor},
  breaklinks=true,
  pdfcreator={LaTeX via pandoc with the Eisvogel template}}
\author{Hugo Roussaffa}
\date{5 août 2026}



%%
%% added
%%


%
% for the background color of the title page
%
\usepackage{pagecolor}
\usepackage{afterpage}
\usepackage[margin=2.5cm,includehead=true,includefoot=true,centering]{geometry}

%
% break urls
%
\PassOptionsToPackage{hyphens}{url}

%
% When using babel or polyglossia with biblatex, loading csquotes is recommended
% to ensure that quoted texts are typeset according to the rules of your main language.
%
\usepackage{csquotes}

%
% captions
%
\definecolor{caption-color}{HTML}{777777}
\usepackage[font={stretch=1.2}, textfont={color=caption-color}, position=top, skip=4mm, labelfont=bf, singlelinecheck=false, justification=raggedright]{caption}
\setcapindent{0em}

%
% blockquote
%
\definecolor{blockquote-border}{RGB}{221,221,221}
\definecolor{blockquote-text}{RGB}{119,119,119}
\usepackage{mdframed}
\newmdenv[rightline=false,bottomline=false,topline=false,linewidth=3pt,linecolor=blockquote-border,skipabove=\parskip]{customblockquote}
\renewenvironment{quote}{\begin{customblockquote}\list{}{\rightmargin=0em\leftmargin=0em}%
\item\relax\color{blockquote-text}\ignorespaces}{\unskip\unskip\endlist\end{customblockquote}}

%
% Source Sans Pro as the default font family
% Source Code Pro for monospace text
%
% 'default' option sets the default
% font family to Source Sans Pro, not \sfdefault.
%
\ifnum 0\ifxetex 1\fi\ifluatex 1\fi=0 % if pdftex
    \usepackage[default]{sourcesanspro}
  \usepackage{sourcecodepro}
  \else % if not pdftex
    \usepackage[default]{sourcesanspro}
  \usepackage{sourcecodepro}

  % XeLaTeX specific adjustments for straight quotes: https://tex.stackexchange.com/a/354887
  % This issue is already fixed (see https://github.com/silkeh/latex-sourcecodepro/pull/5) but the
  % fix is still unreleased.
  % TODO: Remove this workaround when the new version of sourcecodepro is released on CTAN.
  \ifxetex
    \makeatletter
    \defaultfontfeatures[\ttfamily]
      { Numbers   = \sourcecodepro@figurestyle,
        Scale     = \SourceCodePro@scale,
        Extension = .otf }
    \setmonofont
      [ UprightFont    = *-\sourcecodepro@regstyle,
        ItalicFont     = *-\sourcecodepro@regstyle It,
        BoldFont       = *-\sourcecodepro@boldstyle,
        BoldItalicFont = *-\sourcecodepro@boldstyle It ]
      {SourceCodePro}
    \makeatother
  \fi
  \fi

%
% heading color
%
\definecolor{heading-color}{RGB}{40,40,40}
\addtokomafont{section}{\color{heading-color}}
% When using the classes report, scrreprt, book,
% scrbook or memoir, uncomment the following line.
%\addtokomafont{chapter}{\color{heading-color}}

%
% variables for title, author and date
%
\usepackage{titling}
\title{}
\author{Hugo Roussaffa}
\date{5 août 2026}

%
% tables
%

\definecolor{table-row-color}{HTML}{F5F5F5}
\definecolor{table-rule-color}{HTML}{999999}

%\arrayrulecolor{black!40}
\arrayrulecolor{table-rule-color}     % color of \toprule, \midrule, \bottomrule
\setlength\heavyrulewidth{0.3ex}      % thickness of \toprule, \bottomrule
\renewcommand{\arraystretch}{1.3}     % spacing (padding)


%
% remove paragraph indentation
%
\setlength{\parindent}{0pt}
\setlength{\parskip}{6pt plus 2pt minus 1pt}
\setlength{\emergencystretch}{3em}  % prevent overfull lines

%
%
% Listings
%
%


%
% header and footer
%
\usepackage[headsepline,footsepline]{scrlayer-scrpage}

\newpairofpagestyles{eisvogel-header-footer}{
  \clearpairofpagestyles
  \ihead*{}
  \chead*{}
  \ohead*{5 août 2026}
  \ifoot*{Hugo Roussaffa}
  \cfoot*{}
  \ofoot*{\thepage}
  \addtokomafont{pageheadfoot}{\upshape}
}
\pagestyle{eisvogel-header-footer}



%%
%% end added
%%

\begin{document}

%%
%% begin titlepage
%%
\begin{titlepage}
\newgeometry{left=6cm}
\newcommand{\colorRule}[3][black]{\textcolor[HTML]{#1}{\rule{#2}{#3}}}
\begin{flushleft}
\noindent
\\[-1em]
\color[HTML]{5F5F5F}
\makebox[0pt][l]{\colorRule[435488]{1.3\textwidth}{4pt}}
\par
\noindent

{
  \setstretch{1.4}
  \vfill
  \noindent {\huge \textbf{\textsf{}}}
    \vskip 2em
  \noindent {\Large \textsf{Hugo Roussaffa}}
  \vfill
}

\noindent
\includegraphics[width=35mm, left]{../logo\_yapuka.png}

\textsf{5 août 2026}
\end{flushleft}
\end{titlepage}
\restoregeometry
\pagenumbering{arabic}

%%
%% end titlepage
%%



\renewcommand*\contentsname{Sommaire}
{
\hypersetup{linkcolor=}
\setcounter{tocdepth}{3}
\tableofcontents
\newpage
}
\section{Profils utilisateurs
HydroScope}\label{profils-utilisateurs-hydroscope}

\subsection{Profil 1 : Administrateur
plateforme}\label{profil-1-administrateur-plateforme}

\textbf{Description} --- Agent technique (OEIL) responsable de
l'exploitation, de la sécurité et de la maintenance de la plateforme
HydroScope. Gère les comptes utilisateurs, les droits d'accès, les
paramètres système et l'audit des usages.

\textbf{Objectifs} :

\begin{enumerate}
\def\labelenumi{\arabic{enumi}.}
\tightlist
\item
  Sécuriser les accès et garantir la conformité RGPD
\item
  Maintenir la disponibilité et les performances du système
\item
  Paramétrer les profils utilisateurs et les droits d'accès
\item
  Assurer la traçabilité complète des opérations
\item
  Suivre les usages pour adapter la plateforme aux besoins
\end{enumerate}

{\def\LTcaptype{none} % do not increment counter
\begin{longtable}[]{@{}
  >{\raggedright\arraybackslash}p{(\linewidth - 6\tabcolsep) * \real{0.2500}}
  >{\raggedright\arraybackslash}p{(\linewidth - 6\tabcolsep) * \real{0.2500}}
  >{\raggedright\arraybackslash}p{(\linewidth - 6\tabcolsep) * \real{0.2500}}
  >{\raggedright\arraybackslash}p{(\linewidth - 6\tabcolsep) * \real{0.2500}}@{}}
\toprule\noalign{}
\begin{minipage}[b]{\linewidth}\raggedright
\#
\end{minipage} & \begin{minipage}[b]{\linewidth}\raggedright
Cas d'usage
\end{minipage} & \begin{minipage}[b]{\linewidth}\raggedright
US
\end{minipage} & \begin{minipage}[b]{\linewidth}\raggedright
Phrase agile
\end{minipage} \\
\midrule\noalign{}
\endhead
\bottomrule\noalign{}
\endlastfoot
CU-ADMIN-01 & Créer un compte utilisateur et lui attribuer un profil &
US9.1 & En tant qu'administrateur plateforme, je veux créer un compte
utilisateur et lui attribuer un profil afin que chaque acteur puisse
accéder à la plateforme avec les droits appropriés. \\
CU-ADMIN-02 & Configurer les rôles et les droits d'accès par profil &
US9.3 & En tant qu'administrateur plateforme, je veux configurer les
rôles et les droits d'accès par profil afin de garantir que chaque
utilisateur ne dispose que des fonctionnalités nécessaires à son
rôle. \\
CU-ADMIN-03 & Restreindre l'accès à certaines données sensibles selon le
profil & US9.4 & En tant qu'administrateur plateforme, je veux
restreindre l'accès aux données sensibles selon le profil afin de
protéger les informations réglementées. \\
CU-ADMIN-04 & Consulter le journal des actions et détecter des anomalies
& US10.3 & En tant qu'administrateur plateforme, je veux consulter le
journal des actions afin de détecter des comportements anormaux et
assurer la sécurité du système. \\
CU-ADMIN-05 & Lancer un audit d'usage pour adapter la plateforme aux
besoins & US10.6 & En tant qu'administrateur plateforme, je veux lancer
un audit d'usage afin d'identifier les fonctionnalités les plus
utilisées et adapter la plateforme aux besoins réels. \\
\end{longtable}
}

\begin{center}\rule{0.5\linewidth}{0.5pt}\end{center}

\subsection{Profil 2 : Administrateur expert
data}\label{profil-2-administrateur-expert-data}

\textbf{Description} --- Administrateur des données responsable de
l'alimentation, de la qualité, du catalogue, des calculs d'indicateurs
et de la construction de l'aide à la décision. Gère les référentiels
géographiques et méthodologiques.

\textbf{Objectifs} :

\begin{enumerate}
\def\labelenumi{\arabic{enumi}.}
\tightlist
\item
  Alimenter la plateforme avec des données fiables et normalisées
\item
  Qualifier la qualité des données et détecter les anomalies
\item
  Construire et maintenir les référentiels géographiques et
  méthodologiques
\item
  Produire les indicateurs et les rendre accessibles aux autres profils
\item
  Construire l'aide à la décision (AMC, seuils, pondérations) et gérer
  les exports
\end{enumerate}

{\def\LTcaptype{none} % do not increment counter
\begin{longtable}[]{@{}
  >{\raggedright\arraybackslash}p{(\linewidth - 6\tabcolsep) * \real{0.2500}}
  >{\raggedright\arraybackslash}p{(\linewidth - 6\tabcolsep) * \real{0.2500}}
  >{\raggedright\arraybackslash}p{(\linewidth - 6\tabcolsep) * \real{0.2500}}
  >{\raggedright\arraybackslash}p{(\linewidth - 6\tabcolsep) * \real{0.2500}}@{}}
\toprule\noalign{}
\begin{minipage}[b]{\linewidth}\raggedright
\#
\end{minipage} & \begin{minipage}[b]{\linewidth}\raggedright
Cas d'usage
\end{minipage} & \begin{minipage}[b]{\linewidth}\raggedright
US
\end{minipage} & \begin{minipage}[b]{\linewidth}\raggedright
Phrase agile
\end{minipage} \\
\midrule\noalign{}
\endhead
\bottomrule\noalign{}
\endlastfoot
CU-DATA-01 & Importer un nouveau jeu de données et le normaliser &
US1.1, US1.4 & En tant qu'administrateur expert data, je veux importer
un nouveau jeu de données et le normaliser afin que la plateforme
dispose de données fiables et homogènes. \\
CU-DATA-02 & Détecter des valeurs aberrantes et qualifier la fiabilité &
US2.1, US2.5 & En tant qu'administrateur expert data, je veux détecter
des valeurs aberrantes et qualifier la fiabilité afin de garantir la
confiance à accorder aux données produites. \\
CU-DATA-03 & Gérer un bassin versant et ses captages dans le référentiel
& US3.1, US3.2 & En tant qu'administrateur expert data, je veux gérer un
bassin versant et ses captages dans le référentiel afin de maintenir la
cohérence spatiale du système. \\
CU-DATA-04 & Calculer un indicateur et versionner sa méthode & US4.1,
US4.6 & En tant qu'administrateur expert data, je veux calculer un
indicateur et versionner sa méthode afin de garantir la traçabilité et
la reproductibilité des résultats. \\
CU-DATA-05 & Construire un indice composite et documenter les méthodes &
US5.1, US5.6 & En tant qu'administrateur expert data, je veux construire
un indice composite et documenter les méthodes afin de fournir une base
transparente pour l'aide à la décision. \\
\end{longtable}
}

\begin{center}\rule{0.5\linewidth}{0.5pt}\end{center}

\subsection{Profil 3 : Expert métier eau
potable}\label{profil-3-expert-muxe9tier-eau-potable}

\textbf{Description} --- Ingénieur/e spécialisé(e) en eau potable qui
consulte les indicateurs et approfondit l'analyse (AMC, scénarios,
comparaisons). N'effectue pas de gestion de données en back-end mais
exploite les résultats pour son expertise métier.

\textbf{Objectifs} : 1. Conduire l'analyse multicritère pour évaluer des
scénarios 2. Consulter les indicateurs de menace sur l'eau potable d'un
territoire 3. Suivre les tendances et détecter les situations critiques
4. Produire des synthèses pour alimenter la décision 5. Alterner entre
les vues graphiques pour une analyse multi-échelle

{\def\LTcaptype{none} % do not increment counter
\begin{longtable}[]{@{}
  >{\raggedright\arraybackslash}p{(\linewidth - 8\tabcolsep) * \real{0.2000}}
  >{\raggedright\arraybackslash}p{(\linewidth - 8\tabcolsep) * \real{0.2000}}
  >{\raggedright\arraybackslash}p{(\linewidth - 8\tabcolsep) * \real{0.2000}}
  >{\raggedright\arraybackslash}p{(\linewidth - 8\tabcolsep) * \real{0.2000}}
  >{\raggedright\arraybackslash}p{(\linewidth - 8\tabcolsep) * \real{0.2000}}@{}}
\toprule\noalign{}
\begin{minipage}[b]{\linewidth}\raggedright
\#
\end{minipage} & \begin{minipage}[b]{\linewidth}\raggedright
Cas d'usage
\end{minipage} & \begin{minipage}[b]{\linewidth}\raggedright
US
\end{minipage} & \begin{minipage}[b]{\linewidth}\raggedright
Question concrète
\end{minipage} & \begin{minipage}[b]{\linewidth}\raggedright
Phrase agile
\end{minipage} \\
\midrule\noalign{}
\endhead
\bottomrule\noalign{}
\endlastfoot
CU-EXPERT-01 & Conduire l'analyse multicritère (pondérations, scénarios,
comparaisons, contributions) & US5.2, US5.3, US5.4, US5.5 & Sachant que
le BVAEP de la Dumbéa a un statut PPE absent, une surface agricole de 15
\%, 3 ICPE en amont et une géologie ruisselante, quel poids attribuer à
chaque indicateur pour que le classement de criticité reflète le risque
réel de dégradation ? & En tant qu'expert métier eau potable, je veux
conduire l'analyse multicritère et paramétrer les pondérations afin de
créer un score de criticité adapté au contexte de chaque territoire. \\
CU-EXPERT-02 & Consulter les graphiques temporels et alterner les vues &
US6.4, US6.24 & Entre 2021 et 2025, comment a évolué la surface cumulée
d'incendies sur le BVAEP de la Dumbéa, et cette hausse est-elle corrélée
aux surfaces d'érosion identifiées la même année ? & En tant qu'expert
métier eau potable, je veux consulter l'évolution temporelle des
indicateurs et basculer les vues graphiques (changement, répartition,
carte H3) afin d'identifier les dynamiques croisées impactant la
ressource. \\
CU-EXPERT-03 & Identifier des tendances et détecter des dépassements &
US7.3, US7.2 & Le BBR du captage de la Dumbéa est-il passé sous le seuil
de déficit en période d'étiage sur les deux dernières années, indiquant
une tension quantitative sur la ressource ? & En tant qu'expert métier
eau potable, je veux identifier les tendances des indicateurs et être
alerté par les dépassements de seuil que j'ai définis, afin d'anticiper
les périodes de tension sur la ressource. \\
CU-EXPERT-04 & Consulter la fiche synthétique du BVAEP prioritaire &
US7.4, US7.5 & Quels indicateurs contribuent le plus au score de
criticité final du BVAEP de la Dumbéa : est-ce le statut PPE, la densité
d'ICPE, la surface d'incendies cumulés, ou la géologie ? & En tant
qu'expert métier eau potable, je veux consulter la fiche synthétique
d'un territoire prioritaire afin d'identifier rapidement les facteurs de
vulnérabilité dominants avant de recommander une action. \\
CU-EXPERT-05 & Utiliser l'aide à l'interprétation avancée pour prioriser
& US7.6 & Parmi les 5 captages prioritaires du Grand Nouméa, lequel
cumule le plus de facteurs de risque aggravants : absence de PPE + BBR
déficitaire + forte densité d'ICPE + superficie brûlée récurrente ? & En
tant qu'expert métier eau potable, je veux prioriser les territoires
avec l'aide à l'interprétation avancée afin d'établir une hiérarchie
d'actions justifiée partagée avec les décideurs. \\
\end{longtable}
}

\begin{center}\rule{0.5\linewidth}{0.5pt}\end{center}

\subsection{Profil 4 : Décideur
métier}\label{profil-4-duxe9cideur-muxe9tier}

\textbf{Description} --- Élu, directeur de service ou président
d'association disposant d'un pouvoir de décision métier. Utilise
HydroScope pour piloter les orientations stratégiques et opérationnelles
en matière de gestion de l'eau potable.

\textbf{Exemples concrets} : - Commune de référence : Païta (98818),
commune rurale du Grand Nouméa avec captages et zones de protection à
gérer - Captage exemple : Captage de la source de Païta (captage
d'alimentation en eau potable du secteur ouest, soumis à un périmètre de
protection de 1ère catégorie) - Indicateur exemple : Indicateur de
couverture du périmètre de protection (\% de surface du BV effectivement
protégée par un arrêté préfectoral de PPE) - Source de données : Arrêtés
préfectoraux de PPE (textes réglementaires), données cadastrales des
communes (limites communales pour l'intersection BV), référentiel H3 de
la DAVAR - Format de fichier : PDF (rapports d'audit et rapports pour le
conseil municipal), CSV (export de données chiffrées pour tableaux de
bord départementaux) - Technologie : MapLibre (consultation
cartographique en lecture seule), PostGIS (requêtes spatiales pour la
génération de rapports PDF)

\textbf{Objectifs} : 1. Disposer d'une vue synthétique de la qualité de
l'eau sur son territoire 2. Comparer des territoires pour arbitrer des
investissements 3. Suivre les tendances et détecter les situations
critiques 4. Produire des rapports formalisés pour les instances
décisionnelles 5. Consulter l'aide à la décision pour éclairer ses choix

{\def\LTcaptype{none} % do not increment counter
\begin{longtable}[]{@{}
  >{\raggedright\arraybackslash}p{(\linewidth - 6\tabcolsep) * \real{0.2500}}
  >{\raggedright\arraybackslash}p{(\linewidth - 6\tabcolsep) * \real{0.2500}}
  >{\raggedright\arraybackslash}p{(\linewidth - 6\tabcolsep) * \real{0.2500}}
  >{\raggedright\arraybackslash}p{(\linewidth - 6\tabcolsep) * \real{0.2500}}@{}}
\toprule\noalign{}
\begin{minipage}[b]{\linewidth}\raggedright
\#
\end{minipage} & \begin{minipage}[b]{\linewidth}\raggedright
Cas d'usage
\end{minipage} & \begin{minipage}[b]{\linewidth}\raggedright
US
\end{minipage} & \begin{minipage}[b]{\linewidth}\raggedright
Phrase agile
\end{minipage} \\
\midrule\noalign{}
\endhead
\bottomrule\noalign{}
\endlastfoot
CU-DECISION-01 & Consulter un tableau de bord synthétique et une fiche
territoire & US6.6, US6.7 & En tant que décideur métier, je veux
consulter un tableau de bord synthétique et une fiche territoire afin de
disposer d'une vue d'ensemble de la qualité de l'eau sur mon
territoire. \\
CU-DECISION-02 & Comparer deux bassins versants pour prioriser des
investissements & US6.5, US7.5 & En tant que décideur métier, je veux
comparer deux bassins versants afin d'arbitrer les investissements en
fonction des besoins identifiés. \\
CU-DECISION-03 & Identifier une tendance critique et détecter un
dépassement de seuil & US7.3, US7.2 & En tant que décideur métier, je
veux identifier une tendance critique et détecter un dépassement de
seuil afin de réagir rapidement face à une situation urgente. \\
CU-DECISION-04 & Générer un rapport PDF pour un conseil municipal &
US8.3 & En tant que décideur métier, je veux générer un rapport PDF afin
de présenter les résultats au conseil municipal et formaliser la prise
de décision. \\
CU-DECISION-05 & Consulter l'aide à l'interprétation avancée pour
éclairer une décision & US7.6 & En tant que décideur métier, je veux
consulter l'aide à l'interprétation avancée afin de comprendre les
causes sous-jacentes des anomalies détectées. \\
\end{longtable}
}

\begin{center}\rule{0.5\linewidth}{0.5pt}\end{center}

\subsection{Profil 5 : Intéressé
public}\label{profil-5-intuxe9ressuxe9-public}

\textbf{Description} --- Citoyen, habitant ou journaliste souhaitant
consulter la qualité de l'eau sur son territoire. Accès anonyme,
interface simplifiée et filtrable, lecture seule stricte (aucun export,
aucun partage).

\textbf{Exemples concrets} : - Commune de référence : Koné (98811),
chef-lieu de la Province Nord, commune côtière avec captages et BV
d'alimentation en eau potable - Captage exemple : Captage de la rivière
Koné (captage d'alimentation du centre-ville de Koné, soumis à un
périmètre de protection) - Indicateur exemple : Résultat de contrôle
microbiologique (E. coli, coliformes thermotolérants) affiché sous forme
de valeur chiffrée sur survol de la carte - Source de données : Données
publiques de la DAVAR et des Provinces (résultats de prélèvements rendus
publics), référentiel géographique des BV et captages - Format de
fichier : Aucun export disponible (lecture seule stricte) --- les
données sont consultées uniquement via l'interface web - Technologie :
MapLibre (carte interactive en lecture seule), OGC WMS (couches
cartographiques publiques), H3 (indexation spatiale pour le filtrage par
commune)

\textbf{Objectifs} : 1. Vérifier la qualité de l'eau potable sur son
lieu de résidence 2. Explorer la carte des captages et bassins versants
de sa commune 3. Comprendre les indicateurs de qualité grâce à des
données chiffrées 4. Accéder aux périmètres de protection et aux zones à
risque 5. S'informer sur l'état de la ressource en eau sans outil
technique

{\def\LTcaptype{none} % do not increment counter
\begin{longtable}[]{@{}
  >{\raggedright\arraybackslash}p{(\linewidth - 6\tabcolsep) * \real{0.2500}}
  >{\raggedright\arraybackslash}p{(\linewidth - 6\tabcolsep) * \real{0.2500}}
  >{\raggedright\arraybackslash}p{(\linewidth - 6\tabcolsep) * \real{0.2500}}
  >{\raggedright\arraybackslash}p{(\linewidth - 6\tabcolsep) * \real{0.2500}}@{}}
\toprule\noalign{}
\begin{minipage}[b]{\linewidth}\raggedright
\#
\end{minipage} & \begin{minipage}[b]{\linewidth}\raggedright
Cas d'usage
\end{minipage} & \begin{minipage}[b]{\linewidth}\raggedright
US
\end{minipage} & \begin{minipage}[b]{\linewidth}\raggedright
Phrase agile
\end{minipage} \\
\midrule\noalign{}
\endhead
\bottomrule\noalign{}
\endlastfoot
CU-PUBLIC-01 & Vérifier la qualité de l'eau près de chez soi via la
carte & US6.1, US6.2, US6.20 & En tant qu'intéressé public, je veux
vérifier la qualité de l'eau près de chez moi via la carte afin de
connaître l'état de la ressource en eau dans mon quartier. \\
CU-PUBLIC-02 & Rechercher ma commune ou mon captage et consulter les
indicateurs & US6.10, US6.22 & En tant qu'intéressé public, je veux
rechercher ma commune ou mon captage et consulter les indicateurs afin
de trouver rapidement l'information pertinente. \\
CU-PUBLIC-03 & Consulter les couches thématiques et zoomer sur un
territoire & US6.3, US6.21 & En tant qu'intéressé public, je veux
consulter les couches thématiques et zoomer sur un territoire afin
d'explorer visuellement les zones de ma commune. \\
CU-PUBLIC-04 & Appliquer un preset pour voir les territoires les plus
exposés & US6.13 & En tant qu'intéressé public, je veux appliquer un
preset afin de visualiser rapidement les territoires les plus concernés
par les problèmes de qualité de l'eau. \\
CU-PUBLIC-05 & Trier la liste des indicateurs par distance pour trouver
les plus proches & US6.17, US6.25 & En tant qu'intéressé public, je veux
trier la liste des indicateurs par distance afin de trouver les données
les plus proches de mon lieu de résidence. \\
\end{longtable}
}

\end{document}
